\documentclass[10pt, letterpaper]{article}

% Packages:
\usepackage[
    ignoreheadfoot, % set margins without considering header and footer
    top=2 cm, % seperation between body and page edge from the top
    bottom=2 cm, % seperation between body and page edge from the bottom
    left=2 cm, % seperation between body and page edge from the left
    right=2 cm, % seperation between body and page edge from the right
    footskip=1.0 cm, % seperation between body and footer
    % showframe % for debugging
]{geometry} % for adjusting page geometry
\usepackage[explicit]{titlesec} % for customizing section titles
\usepackage{tabularx} % for making tables with fixed width columns
\usepackage{array} % tabularx requires this
\usepackage[dvipsnames]{xcolor} % for coloring text
\definecolor{primaryColor}{RGB}{0, 79, 144} % define primary color
\definecolor{accentColor}{RGB}{224, 82, 99}
\definecolor{tealColor}{RGB}{0, 137, 123}
\definecolor{purpleColor}{RGB}{112, 72, 166}
\definecolor{goldColor}{RGB}{218, 146, 35}
\definecolor{softGray}{RGB}{95, 105, 120}
\usepackage{enumitem} % for customizing lists
\usepackage{fontawesome5} % for using icons
\usepackage{amsmath} % for math
\usepackage[UTF8, scheme=plain]{ctex} % for Chinese support
\usepackage[
    pdftitle={[你的姓名]的简历},
    pdfauthor={[你的姓名]},
    pdfcreator={LaTeX with RenderCV},
    colorlinks=true,
    linkcolor=purpleColor,
    urlcolor=accentColor
]{hyperref} % for links, metadata and bookmarks
\usepackage[pscoord]{eso-pic} % for floating text on the page
\usepackage{calc} % for calculating lengths
\usepackage{bookmark} % for bookmarks
\usepackage{lastpage} % for getting the total number of pages
\usepackage{changepage} % for one column entries (adjustwidth environment)
\usepackage{paracol} % for two and three column entries
\usepackage{ifthen} % for conditional statements
\usepackage{needspace} % for avoiding page brake right after the section title
\usepackage{iftex} % check if engine is pdflatex, xetex or luatex

% Ensure that generate pdf is machine readable/ATS parsable:
\ifPDFTeX
    \input{glyphtounicode}
    \pdfgentounicode=1
    \usepackage[T1]{fontenc}
    \usepackage[utf8]{inputenc}
    \usepackage{lmodern}
\fi

\usepackage[default, type1]{sourcesanspro}

% Some settings:
\AtBeginEnvironment{adjustwidth}{\partopsep0pt} % remove space before adjustwidth environment
\pagestyle{empty} % no header or footer
\setcounter{secnumdepth}{0} % no section numbering
\setlength{\parindent}{0pt} % no indentation
\setlength{\topskip}{0pt} % no top skip
\setlength{\columnsep}{0.15cm} % set column seperation
\makeatletter
\let\ps@customFooterStyle\ps@plain % Copy the plain style to customFooterStyle
\patchcmd{\ps@customFooterStyle}{\thepage}{
    \color{softGray}\textit{\small [你的姓名] - 第 \textcolor{accentColor}{\thepage{}} 页，共 \pageref*{LastPage} 页}
}{}{} % replace number by desired string
\makeatother
\pagestyle{customFooterStyle}

\titleformat{\section}{
    % avoid page braking right after the section title
    \needspace{4\baselineskip}
    % make the font size of the section title large and color it with the primary color
    \Large\color{primaryColor}
}{}{}{
    % print bold title, give 0.15 cm space and draw a line of 0.8 pt thickness
    % from the end of the title to the end of the body
    \textcolor{primaryColor}{\textbf{#1}}\hspace{0.15cm}{\color{accentColor}\titlerule[0.8pt]}\hspace{-0.1cm}
}[] % section title formatting

\titlespacing{\section}{
    % left space:
    -1pt
}{
    % top space:
    0.3 cm
}{
    % bottom space:
    0.2 cm
} % section title spacing

\renewcommand\labelitemi{\textcolor{tealColor}{\footnotesize\faAngleRight}} % custom bullet points
\renewcommand\labelitemii{\textcolor{goldColor}{\scriptsize\faCircle}} % custom nested bullet points
\newenvironment{highlights}{
    \begin{itemize}[
        topsep=0.10 cm,
        parsep=0.10 cm,
        partopsep=0pt,
        itemsep=0pt,
        leftmargin=0.4 cm + 10pt
    ]
}{\end{itemize}} % new environment for highlights

\newenvironment{highlightsforbulletentries}{
    \begin{itemize}[
        topsep=0.10 cm,
        parsep=0.10 cm,
        partopsep=0pt,
        itemsep=0pt,
        leftmargin=10pt
    ]
}{\end{itemize}} % new environment for highlights for bullet entries


\newenvironment{onecolentry}{
    \begin{adjustwidth}{
        0.2 cm + 0.00001 cm
    }{
        0.2 cm + 0.00001 cm
    }
}{\end{adjustwidth}} % new environment for one column entries

\newenvironment{twocolentry}[2][]{\onecolentry
    \def\secondColumn{#2}
    \setcolumnwidth{\fill, 4.5 cm}
    \begin{paracol}{2}
}{\switchcolumn \raggedleft \secondColumn
    \end{paracol}
    \endonecolentry} % new environment for two column entries

\newenvironment{threecolentry}[3][]{\onecolentry
    \def\thirdColumn{#3}
    \setcolumnwidth{1.2 cm, \fill, 4.5 cm}
    \begin{paracol}{3}
    {\raggedright #2} \switchcolumn
}{\switchcolumn \raggedleft \thirdColumn
    \end{paracol}
    \endonecolentry} % new environment for three column entries

\newenvironment{header}{
    \setlength{\topsep}{0pt}\par\kern\topsep\centering\color{primaryColor}\linespread{1.5}
}{\par\kern\topsep} % new environment for the header

\newcommand{\placelastupdatedtext}{% \placetextbox{<horizontal pos>}{<vertical pos>}{<stuff>}
  \AddToShipoutPictureFG*{% Add <stuff> to current page foreground
    \put(
        \LenToUnit{\paperwidth-2 cm-0.2 cm+0.05cm},
        \LenToUnit{\paperheight-1.0 cm}
    ){\vtop{{\null}\makebox[0pt][c]{
        \small\color{gray}\textit{最后更新于 2026 年 4 月}\hspace{\widthof{最后更新于 2026 年 4 月}}
    }}}%
  }%
}%

% save the original href command in a new command:
\let\hrefWithoutArrow\href

% new command for external links:
\renewcommand{\href}[2]{\hrefWithoutArrow{#1}{\ifthenelse{\equal{#2}{}}{ }{#2 }\raisebox{.15ex}{\footnotesize \faExternalLink*}}}


\pagecolor{white}
\begin{document}
    \newcommand{\AND}{\unskip
        \cleaders\copy\ANDbox\hskip\wd\ANDbox
        \ignorespaces
    }
    \newsavebox\ANDbox
    \sbox\ANDbox{}

    \placelastupdatedtext
    \begin{header}
        \fontsize{30 pt}{30 pt}
        \textbf{[你的姓名]}

        \vspace{0.3 cm}

        \normalsize
        \mbox{\textcolor{accentColor}{\footnotesize\faMapMarker*}\hspace*{0.13cm}[你的城市]}%
        \kern 0.25 cm%
        \AND%
        \kern 0.25 cm%
        \mbox{\hrefWithoutArrow{mailto:[你的邮箱]}{\textcolor{tealColor}{\footnotesize\faEnvelope[regular]}\hspace*{0.13cm}[你的邮箱]}}%
        \kern 0.25 cm%
        \AND%
        \kern 0.25 cm%
        \mbox{\hrefWithoutArrow{tel:[你的电话]}{\textcolor{goldColor}{\footnotesize\faPhone*}\hspace*{0.13cm}[你的电话]}}%
        \kern 0.25 cm%
        \AND%
        \kern 0.25 cm%
        \mbox{\hrefWithoutArrow{https://github.com/[你的GitHub用户名]}{\textcolor{purpleColor}{\footnotesize\faGithub}\hspace*{0.13cm}github.com/[你的GitHub用户名]}}%
    \end{header}

    \vspace{0.3 cm - 0.3 cm}


    \section{教育背景}

        \begin{threecolentry}{\textbf{本科}}{
            2023.09 -- 2027.06（预计）
        }
            \textbf{[你的大学名称]}，[你的专业名称]
            \begin{highlights}
                \item \textbf{GPA:} [X.X/4.0] | 排名: \textbf{[X/N（前 X\%）]} % 可选
                \item \textbf{英语水平:} CET-4 \textbf{[成绩]}, CET-6 \textbf{[成绩]} % 可选
                \item \textbf{核心课程:}\par\vspace{0.05cm}
                \noindent
                \begin{minipage}[t]{0.47\linewidth}
                    \textbf{计算机 / 人工智能课程:}
                    \begin{itemize}[topsep=0.05cm, parsep=0pt, itemsep=0pt, leftmargin=1.1em]
                        \item 自然语言处理 \hfill \textbf{A}
                        \item 机器学习导论 \hfill \textbf{A}
                        \item 人工智能 I \hfill \textbf{A}
                        \item 计算机编程 \hfill \textbf{A+}
                        \item 信息科学与技术导论 \hfill \textbf{A}
                    \end{itemize}
                \end{minipage}\hfill
                \begin{minipage}[t]{0.47\linewidth}
                    \textbf{数学基础:}
                    \begin{itemize}[topsep=0.05cm, parsep=0pt, itemsep=0pt, leftmargin=1.1em]
                        \item 凸优化 \hfill \textbf{A+}
                        \item 概率论与数理统计 \hfill \textbf{A}
                        \item 线性代数与应用 \hfill \textbf{A}
                        \item 微积分 I \hfill \textbf{A}
                        \item 微积分 II \hfill \textbf{A-}
                    \end{itemize}
                \end{minipage}
            \end{highlights}
        \end{threecolentry}



    \section{科研经历}

        \begin{twocolentry}{
            [你的研究地点]

        [开始日期] -- [结束日期]
        }
            \textbf{[你的研究项目标题]} \\
            {\small\textcolor{softGray}{指导教师: [指导教师姓名]}} \hfill \textit{[作者身份], [发表情况]}
            \begin{highlights}
                \item \textbf{背景:} [请在此处描述研究背景和问题陈述]
                \item \textbf{方法:} [请在此处描述你的研究方法和技术方案]
                \begin{itemize}[topsep=0.05cm, parsep=0pt, itemsep=0pt, leftmargin=*]
                    \item {[方法要点 1]}
                    \item {[方法要点 2]}
                    \item {[方法要点 3]}
                \end{itemize}
                \item \textbf{贡献:}
                \begin{itemize}[topsep=0.05cm, parsep=0pt, itemsep=0pt, leftmargin=*]
                    \item {[研究贡献 1]}
                    \item {[研究贡献 2]}
                    \item {[研究贡献 3]}
                \end{itemize}
            \end{highlights}
        \end{twocolentry}

        \vspace{0.2 cm}

        \begin{twocolentry}{
            [你的研究地点]

        [开始日期] -- [结束日期]
        }
            \textbf{[你的研究项目标题]} \\
            {\small\textcolor{softGray}{指导教师: [指导教师姓名]}} \hfill \textit{[作者身份]}
            \begin{highlights}
                \item \textbf{背景:} [请在此处描述研究背景和问题陈述]
                \item \textbf{方法:} [请在此处描述你的研究方法和技术方案]
                \begin{itemize}[topsep=0.05cm, parsep=0pt, itemsep=0pt, leftmargin=*]
                    \item {[方法要点 1]}
                    \item {[方法要点 2]}
                \end{itemize}
                \item \textbf{贡献:}
                \begin{itemize}[topsep=0.05cm, parsep=0pt, itemsep=0pt, leftmargin=*]
                    \item {[研究贡献 1]}
                    \item {[研究贡献 2]}
                    \item {[研究贡献 3]}
                \end{itemize}
            \end{highlights}
        \end{twocolentry}

        \vspace{0.2 cm}

        \begin{twocolentry}{
            [你的研究地点]

        [开始日期] -- [结束日期]
        }
            \textbf{[你的项目标题]} \\
            {\small\textcolor{softGray}{[实验室 / 组织名称]，指导教师: [指导教师姓名]}}
            \begin{highlights}
                \item \textbf{工作:} [请在此处描述你的项目工作内容和贡献]
            \end{highlights}
        \end{twocolentry}


    \section{项目经历}

        \begin{twocolentry}{
            \href{[项目网址]}{[项目网址]}
        }
            \textbf{[项目名称]} \\
            {\small\textcolor{softGray}{[项目描述]}}
            \begin{highlights}
                \item \textbf{工作:} [请在此处描述你的项目工作内容和贡献]
            \end{highlights}
        \end{twocolentry}


    \section{荣誉奖项与教学经历}

        \begin{onecolentry}
            \begin{highlightsforbulletentries}
                \item \textbf{[奖项或荣誉名称]}, [年份] % 可选
                \item \textbf{[课程名称]} 助教，[大学名称]，[学期 / 年份] % 可选
                \item {[职位 / 角色]}，[社团 / 组织名称]，[年份] % 可选
                \item {[其他荣誉或经历]} % 可选
            \end{highlightsforbulletentries}
        \end{onecolentry}



    \section{专业技能}

        \begin{onecolentry}
            \textbf{编程语言与工具:} Python (PyTorch, HuggingFace Transformers), C/C++, Linux, Git, LaTeX
        \end{onecolentry}

        \vspace{0.2 cm}

        \begin{onecolentry}
            \textbf{研究兴趣:} NLP, LLM Reasoning (Chain-of-Thought), Transformer Architecture
        \end{onecolentry}


\end{document}
